\chapter{Glossary}

\begin{description}
  \item[APFS] Apple File System.
  \item[Checkpoint] A transaction state used to locate coherent APFS metadata.
  \item[XID] Transaction identifier. Raw scans must avoid mixing incompatible
  transaction contexts.
  \item[OID] Object identifier. A logical object id is not by itself a safe
  parsed-content cache key.
  \item[OMAP] Object map. The B-tree used to resolve APFS virtual objects to
  physical locations in the correct domain and transaction context.
  \item[FS tree] File-system B-tree containing records used to reconstruct
  namespace and file metadata.
  \item[DIR\_REC] Directory record used for directory membership and names.
  \item[INODE] Inode record used for file identity and metadata.
  \item[Dstream] Data stream metadata that carries logical size information for
  ordinary file content.
  \item[Oracle] A feature-specific validation source. There is no single oracle
  for all APFS product semantics.
  \item[Raw-readable] Bytes can be opened and sampled from a source.
  \item[Raw-supported] Raw output is validated for the requested product mode and
  source class.
  \item[Logical size] The v1 size metric, corresponding to the user-visible file
  length rather than allocated/shared storage attribution.
  \item[Snapshot-retained bytes] Storage retained by APFS snapshot semantics.
  This is not part of raw v1 logical-size output.
  \item[Firmlink] macOS presentation mechanism that contributes to merged
  System/Data namespace behavior.
  \item[NX / NXSB] APFS container superblock (\texttt{nx\_superblock\_t}). The
  block-zero copy is the locator; the authoritative selected NXSB lives in the
  checkpoint descriptor area at the highest valid checkpoint XID.
  \item[APSB] APFS volume superblock. Decoded by the Rust crate under the v1
  feature allowlist; encryption, sealed volumes, and unknown incompatible
  features force \texttt{Unsupported(reason)}.
  \item[Checkpoint map] The per-checkpoint
  \texttt{checkpoint\_map\_phys\_t} ring that maps ephemeral object IDs to
  physical addresses for the selected transaction; the last block in the ring
  must carry \texttt{CHECKPOINT\_MAP\_LAST}.
  \item[Scan XID] The single chosen transaction identifier the parser pins for
  the duration of a run; all object resolution must be performed against this
  scan XID.
  \item[Xfield] The variable-length extended-field area in inode and
  directory-record bodies. SR-015 fixes one source-backed cursor rule:
  values start immediately after the metadata table and each value advances by
  \texttt{round\_up(x\_size, 8)}; \texttt{xf\_used\_data} is the padded
  value-area byte count.
  \item[Sibling link / sibling map] The hard-link records
  (\texttt{j\_sibling\_link\_val\_t}, \texttt{j\_sibling\_map\_val\_t}) that
  unify multiple paths sharing a single file identity.
  \item[decmpfs] The Apple compression xattr family used by the
  HFS-compression mechanism on APFS. SR-009 and SR-017 specify when its
  uncompressed size takes precedence over \texttt{j\_dstream\_t.size}.
  \item[Sparse bytes] The \texttt{INO\_EXT\_TYPE\_SPARSE\_BYTES} xfield that
  records bytes saved by sparse-file allocation; relevant to the size
  precedence table when the dstream size includes sparse holes.
  \item[Fail-closed] The behavioral discipline that turns any unrecognized,
  malformed, or out-of-allowlist condition into a typed hard stop rather than
  a best-effort guess; centralized in
  \texttt{object::validate\_object\_block} for object headers and in SR-016
  for record bodies.
\end{description}
